\section{Performance dimension: Intra-node (node-level performance)}

\marginpar{20\% memory ist ok anscheinend.

Bei MPI aich net so strikt schliesslich braucht Tool selber auch was}

On an individual node, our first goal has to be that we use all the hardware threads efficiently.
The term intra-node hence means all those cores sharing one memory.
Without employing a tool, my method of choice to argue about node-level performance is strong scaling.
On a chip with 128 cores, it makes sense to understand how much we gained by switching from 1 to 2 to 4 to \ldots cores.
It tells us if we have made optimal use of the whole chip.


\begin{assessmentcrime}
 We start an intra-node assessment with a code which does not use most of the memory.
\end{assessmentcrime}

\noindent
We may assume that a node and its memory are somehow balanced.
If we don't use a reasonable amount of this memory for a test run, then there is simply not enough work to do to utilise all cores.
We can barely blame the code.


\subsection{Is there a problem?}

To make an initial assessment how well we utilise the cores, we rely on strong scaling.
Amdahl's law reads

\begin{equation}
 \label{equation:performance-assessment-intra-node}
 t(p_{\text{core}}) = (1-f_{\text{intra-node}}) \cdot t(1)  + f_{\text{intra-node}} t(1)/p_{\text{core}},
\end{equation}
 
\noindent
where $t(p_{\text{core}})$ is the time used on $p_{\text{core}}$ compute units (cores), and $f_{\text{intra-node}}$ is the serial fraction within a code.

\begin{performancemetric}
 If $f_{\text{intra-node}}$ in (\ref{equation:performance-assessment-intra-node}) falls below 80\%, our code suffers from insufficient strong scaling.
 It fails to utilise all cores on the system.
\end{performancemetric}

\noindent
Computing the $f_{\text{intra-node}}$ should be relatively straightforward.
We gradually reduce the core count from $p_{\text{intra-node}}$ to $p_{\text{intra-node}}/2,p_{\text{intra-node}}/4,\ldots$ and plot the timings. 
In OpenMP, we can simply reduce \texttt{OMP\_NUM\_THREADS} to achieve this.
Other tasking frameworks might require us to change the thread count in the Slurm script, or to set some other OS parameters. 


There are two unknowns of interest in (\ref{equation:performance-assessment-intra-node}):
$f_{\text{intra-node}}$ and $t(1)$.
The latter is not really of interest to use for the metric, but we need a reference timing (we could have picked an arbitrary other value).
High school maths tells us that we therefore need only two measurements to determine the two quantities:
Measure twice, insert the results into each other and you have $f_{\text{intra-node}}$ and $t(1)$.
As we want to get our results as quickly as possible, it makes sense to start with the biggest core count and then to assess half the core count.


In practice, I recommend to start with at least three values and hence to determine two $f_{\text{intra-node}}$.
If they differ by magnitudes or are both way below 80\%, I continue reducing the core count and measuring.
In cases where $f_{\text{intra-node}}$ drops suddenly over $p_{\text{intra-node}}$, i.e.~is really depending on the core count, the particular setups around the drop deserve further attention.


If we have a performance flaw, we continue investigating.
